\begin{proposition}[Root-atlas base change and relative duality]\hypertarget{tag-033I}{}\label{push:src:E.3.1}
\StatementRecord{push:src:E.3.1}{proved}{033I}{the normalized torus root charts and the periodic triangular cusp charts; \cite[Chapter III, Theorem 4.12(ii), pp.~134--135, and Theorem 4.10, pp.~132--133]{BS76}; \cite[introduction, equation (1), p.~261; Section 1, pp.~264--265; Proposition 2.1, p.~266]{RRV71}}

Let $k,l$ be positive integers with $\gcd(3k,4l)=1$.
For the root lift $\bar f:X\to\mathfrak B_{k,l}=\mathbb P(3k,4l)$
of the construction in \cref{fam:src:B.3.1}, the direct images
$R^q\bar f_*\mathcal O_X$, for $q=0,1,2$, are locally free of
ranks $1,2,1$. The canonical geometric-point maps
\[
 (R^q\bar f_*\mathcal O_X)\otimes k(b)
 \longrightarrow H^q(X_b,\mathcal O_{X_b})
\]
are isomorphisms, and the trace induces the canonical isomorphism
\[
 (R^2\bar f_*\mathcal O_X)^\vee
 \simeq\bar f_*\omega_{X/\mathfrak B_{k,l}}.
\]
Both comparisons commute with the analytic root-atlas groupoid.
Here the geometric fibre at a finite root point is the smooth torus
on the normalized root chart.
\end{proposition}

\begin{proof}
The fibre-divisor identities $f^*p_1=3kS_1$ and $f^*p_2=4lS_2$
define the root lift. On a full root chart its atlas pullback is
\[
 [J'_j/\mu_{m_j}]
 \mathbin{\times_{[\Delta_{s_j}/\mu_{m_j}]}}
 \Delta_{s_j}=J'_j,
 \qquad (m_1,m_2)=(3k,4l).
\]
The affine quotient action is free and the family $J'_j\to\Delta_{s_j}$
is the smooth proper torus family. The atlas map is \'etale.
At the cusp the local charts are $q=z_0z_1z_2$, with the
two-factor and smooth charts on the other strata. They have smooth
total space and pure two-dimensional Cartier fibres. A base parameter
is a non-zero-divisor, so the atlas family is flat. Properness is
properness of the compact torus and cusp fillings. These local
comparisons make the atlas families proper flat Gorenstein morphisms
of relative dimension two.

A smooth two-torus has structure-sheaf dimensions $(1,2,1)$.
Let $\nu:P\to W$ be the normalization, $P=\mathrm{dP}_6$, let
$C_h,C_v,C_d$ be the three double-curve normalizations, and let
$P_+,P_-$ be the two triple points. Each $C_e$ is $\mathbb P^1$.
Write $\rho_e^0,\rho_e^1:C_e\to P$ for its two paired boundary maps
and $j_e=\nu\rho_e^0=\nu\rho_e^1$. The sheaf sequence is
\begin{equation}\hypertarget{tag-01OY}{}\label{fam:eq:ab-cusp-resolution}
 0\longrightarrow\mathcal O_W\longrightarrow\nu_*\mathcal O_P
 \xrightarrow{\delta^0}\bigoplus_{e=h,v,d}j_{e*}\mathcal O_{C_e}
 \xrightarrow{\delta^1}\mathbb C_{P_+}\oplus\mathbb C_{P_-}
 \longrightarrow0,
\end{equation}
The differentials in this complex are
\begin{equation}\hypertarget{tag-01OJ}{}\label{eq:P5normalization}
\begin{aligned}
 \delta^0(s)_e&=(\rho_e^1)^*s-(\rho_e^0)^*s,\\
 \delta^1(t)&=\bigl((t_h-t_v+t_d)(P_+),
                        (t_h-t_v+t_d)(P_-)\bigr).
\end{aligned}
\end{equation}
Its exactness is local: for $xyz=0$ it is the alternating restriction
complex of the three coordinate planes, their pairwise lines and their
common point. Decomposing convergent power series by the variables
occurring in each monomial gives the augmented simplex complex, which
is exact. The one- and two-branch cases are the same calculation with
fewer planes. This verifies exactness for the boundary maps,
including the torus multipliers in their gluing.

All positive-degree structure-sheaf cohomology of $P,C_e$ vanishes.
Consequently \eqref{fam:eq:ab-cusp-resolution} computes cohomology by
\[
 \mathbb C\xrightarrow{0}\mathbb C^3
 \xrightarrow{\left(\begin{smallmatrix}1&-1&1\\1&-1&1\end{smallmatrix}\right)}
 \mathbb C^2.
\]
It gives $(h^0,h^1,h^2)(\mathcal O_W)=(1,2,1)$.
We also need its multiplication. The branch-intersection resolution
is the augmented semisimplicial sheaf algebra with its restriction face
maps; its local exactness makes the augmentation a multiplicative
quasi-isomorphism. Its Alexander--Whitney product therefore computes
the sheaf-cohomology cup product. In the ordered triangles
$[0,e_1,e_2]$ and $[e_1,e_2,e_1+e_2]$ it is
\[
 \alpha\smile\beta=(\alpha_h\beta_d,\alpha_d\beta_h),
 \qquad [\alpha]\smile[\beta]
 =\alpha_h\beta_d-\alpha_d\beta_h
 \quad\hbox{in }\mathbb C^2/\mathbb C(1,1).
\]
The cocycles $(1,1,0)$ and $(0,1,1)$ pair to one. Thus
$\bigwedge^2H^1(\mathcal O_W)\to H^2(\mathcal O_W)$ is an
isomorphism.
Thus the three dimensions are constant on every
atlas chart, including its centre.

The constant-cohomology theorem
\cite[Chapter III, Theorem 4.12(ii), pp.~134--135]{BS76}
now gives local freeness, and the canonical analytic base-change
maps are those of
\cite[Chapter III, Theorem 4.10, pp.~132--133]{BS76}.
The hypotheses are properness, flatness of the structure sheaf,
a reduced smooth base, and the constant fibre dimensions just
computed. The constant-dimension formulation also appears in
\cite[Section 7, Satz 5, pp.~63--64]{Grauert60}, with
\cite[pp.~35--36]{Grauert1963ErratumD}.
On a double atlas overlap, both pullbacks identify with the direct
image of the same pulled-back family. Functoriality makes these
identifications satisfy the triple-overlap cocycle condition, so
they descend with their stabilizer representations.

For an atlas family $p:Y\to U$, proper analytic duality gives
\[
 Rp_*\omega_{Y/U}[2]
 \simeq R\mathcal Hom_U(Rp_*\mathcal O_Y,\mathcal O_U).
\]
This is the canonical trace identity of
\cite[introduction, equation (1), p.~261; Section 1, pp.~264--265;
Proposition 2.1, p.~266]{RRV71}.
The smooth-space dualizing complex is $K_Y[3]$ and the relative
Gorenstein complex is $\omega_{Y/U}[2]$; see
\cite[Section 3, Theorem 1, p.~81; Proposition 2 and the
complete-intersection calculation, p.~91]{RR70push}.
The locally free cohomology sheaves already obtained have no
positive sheaf Ext against $\mathcal O_U$. Degree $-2$ of the
trace identity is therefore
$p_*\omega_{Y/U}\simeq(R^2p_*\mathcal O_Y)^\vee$.
The trace commutes with the local biholomorphisms on the atlas
and has the same cocycle compatibility. It consequently descends
to the asserted isomorphism on $\mathfrak B_{k,l}$.
\end{proof}